% Unidad U006F. Inventario operatorio completo del Topological Phase Kernel.

\chapter{Inventario operatorio completo y tubería causal del Topological Phase Kernel}
\label{chap:inventario-operatorio-completo-tpk}

El Topological Phase Kernel no es una reducción módulo tres ni un reloj de
ciento ocho posiciones. Es una categoría de estados enriquecidos equipada
con operadores de selección, transporte, lectura, memoria, periodicidad,
prolongación y publicación. Para impedir que una abreviatura sustituya al
objeto, cada operador se registra por su tipo y por la información que
conserva.

\section{Ficha tipada de un operador}

\begin{definition}[Ficha operatoria]
\label{def:u006f-ficha-operatoria}
Una realización canónica del TPK es una séxtupla
\begin{equation}
 \mathcal O=
 (D_{\mathcal O},C_{\mathcal O},f_{\mathcal O},
 I_{\mathcal O},L_{\mathcal O},M_{\mathcal O}),
 \label{eq:u006f-ficha}
\end{equation}
donde:
\begin{enumerate}
 \item \(f_{\mathcal O}:D_{\mathcal O}\to C_{\mathcal O}\) es su regla;
 \item \(I_{\mathcal O}\) es la familia de relaciones conservadas;
 \item \(L_{\mathcal O}\) declara la información publicada o perdida;
 \item \(M_{\mathcal O}\) enumera la memoria que permanece en el estado
 enriquecido.
\end{enumerate}
\end{definition}

Dos realizaciones representan el mismo operador sólo cuando existe una
equivalencia tipada que entrelaza reglas e invariantes. Compartir periodo,
cardinal o salida visible no basta.

\section{Ocho clases estructurales}

El inventario se particiona en:
\begin{enumerate}
 \item[\textbf{C1.}] soportes y estados;
 \item[\textbf{C2.}] selección interna;
 \item[\textbf{C3.}] transporte;
 \item[\textbf{C4.}] lectores y cocientes;
 \item[\textbf{C5.}] memoria y frontera;
 \item[\textbf{C6.}] periodicidad y retorno;
 \item[\textbf{C7.}] prolongación;
 \item[\textbf{C8.}] salidas.
\end{enumerate}
La clase se determina por dominio, codominio y función, no por el nombre de
una cifra asociada.

\section{Catorce agrupaciones funcionales}

\begingroup
\small
\renewcommand{\arraystretch}{1.12}
\begin{longtable}{@{}p{.20\textwidth}p{.29\textwidth}p{.43\textwidth}@{}}
\toprule
Agrupación&Operadores o espacios&Función e invariantes\\
\midrule
\endfirsthead
\toprule
Agrupación&Operadores o espacios&Función e invariantes\\
\midrule
\endhead
Estado y memoria de ruta&
\(\mathcal X_{\rm TPK}^{\rm enr}\), \(\mathcal F_q\)&
Conservan posición, hoja, orientación, residuo, cociente, acarreo, firma,
frontera, cilindro, prefijo, procedencia y memoria.\\

Selección interna&
\(M_{\rm ph}:\{1,\ldots,9\}\to\{-1,0,+1\}\)&
Activa la evaluación aditiva, multiplicativa o el paso de umbral; no mueve
el estado.\\

Transporte espacial&
\(T_N,T_E,T_S,T_O,F_{\rm loc},F_{\rm obs}\)&
Realiza traslaciones cardinales, realimentación de hoja y cronología de dos
cursores.\\

Lectores modulares&
\(\rho_9,q_9,\rho_3^{\rm or},c_3,q_{1000},r_{1000}\)&
Publican fase, orientación y bloque decimal conservando los cocientes de
levantamiento.\\

Bloque y memoria&
\(\mathcal K_{27}=F_{\rm obs}^{27}\), \(\mathcal N_{27}\)&
El primero compone veintisiete transportes; el segundo filtra eventos de
memoria. Igual índice no implica igual operador.\\

Estado dodecafásico&
\(\mathcal R_{12},C_{12},\Theta_{108},\Gamma_{108}\)&
Distingue registro enriquecido, sector, ley cíclica y soporte ordenado.\\

Periodicidad y profundidades&
\(C_{12}\hookrightarrow C_{108}\twoheadrightarrow C_9\),
\(w_{108},w_{1080}\)&
Conserva la extensión no escindida, las palabras de distintas longitudes y
la memoria vertical.\\

Canales \(90/120\)&
\(\mathcal F_6,\mathcal F_8,\mathcal I_{6\to8},
S_{90},S_{120}\)&
Separa incidencias finitas, colas analíticas y sus normalizaciones.\\

Bidireccionalidad&
\(P_\pm,\operatorname{rev},\operatorname{Ext}_\pm,
N_\pm,\beta,C_M\)&
Realiza intercambio, inversión de ruta, prolongación futura y pasada,
valencia, balance y conjugación.\\

Emisión y coinducción&
\(L_0,L_1,R_{36},G_9,\varprojlim\operatorname{Cyl}_m\)&
Eleva la semilla, abre la frontera y organiza prolongaciones finitas y
compatibles.\\

Lectores numéricos&
\(\mathscr C_\pi,\mathscr P_e,F_{\rm av}\)&
Construyen horquillas racionales, propagación factorial y autoescala de
Fibonacci.\\

Proyección excepcional&
\(\Gamma_{A_W}\), peso, soporte proyectivo&
Transporta la misma emisión a incidencia ternaria sin seleccionar cifras
arquimedianas.\\

Atlas y realización&
firma de ruta, atlas \(81/56/13/324\), \(C_M\)&
Clasifica historias y caracteres antes de aplicar una carta dimensional.\\

Validación finita&
censos, ablaciones, hashes y pruebas de independencia&
Comprueba cobertura en la profundidad ejecutada y separa generador,
verificador y datos de contraste.\\
\bottomrule
\end{longtable}
\endgroup

Las ocho clases describen naturaleza matemática; las catorce agrupaciones,
función de ejecución. Son dos gradaciones del mismo inventario.

\section{Registro nominal de treinta y cuatro entradas}

\begin{equation}
 \mathfrak O_{\rm TPK}:=(\mathfrak o_1,\ldots,\mathfrak o_{34}).
 \label{eq:u006f-registro-operadores}
\end{equation}

\begingroup
\scriptsize
\renewcommand{\arraystretch}{1.12}
\begin{longtable}{@{}r p{.19\textwidth}p{.28\textwidth}p{.42\textwidth}@{}}
\toprule
\#&Símbolo&Dominio y codominio&Regla o función\\
\midrule
\endfirsthead
\toprule
\#&Símbolo&Dominio y codominio&Regla o función\\
\midrule
\endhead
1&\(\mathcal A_9\)&\((\mathbb Z/9)^2\), matriz y ley interna&
Célula aritmética aditivo--multiplicativa.\\
2&\((\rho_9,q_9)\)&\(\mathbb Z_{>0}\to
\{1,\ldots,9\}\times\mathbb Z_{\geq0}\)&
División nonádica con representante positivo y cociente.\\
3&\(\rho_3^{\rm or},c_3\)&\(\mathbb Z\to
\{-1,0,+1\}\times\mathbb Z\)&
Levantamiento ternario orientado \(n=\rho_3^{\rm or}(n)+3c_3(n)\).\\
4&\(\iota_{1000}:=(q_{1000},r_{1000})\)&\(\mathbb Z\to
\mathbb Z\times\{0,\ldots,999\}\)&
División euclídea \(N=1000q_{1000}+r_{1000}\).\\
5&\(T_d\)&\(X_9\to X_9\), \(d\in\{N,E,S,O\}\)&
Traslación cardinal orientada.\\
6&\(F_{\rm loc}\)&\(\mathcal Q_{\rm loc}\times
\mathscr D_{\rm card}\to\mathcal Q_{\rm loc}\)&
Realimentación local de posición y hoja.\\
7&\(T_{\min}\)&\(\Omega_{\rm seed}\to\mathcal U_6\)&
Emisor mínimo de los dos cursores, incluida la firma coordenada.\\
8&\(G\)&\(\operatorname{im}s_\Sigma\times
\operatorname{im}s_\Pi\to\{0,\ldots,999\}^6\)&
Acoplamiento decimal coordenada a coordenada.\\
9&\(\Psi_6\)&\(\mathcal U_6\to\mathbb F_3^6\)&
Reducción ternaria del catálogo.\\
10&\(\mathcal R_3\)&\(\mathcal Q_{\rm loc}\times
\mathscr D_{\rm card}^3\to\{-1,0,+1\}\)&
Factor observable local de tres pasos.\\
11&\(\mathcal K_{27}\)&\(\mathcal Q_{\rm obs}\to\mathcal Q_{\rm obs}\)&
Composición \(F_{\rm obs}^{27}\).\\
12&\(\mathcal X_{\rm TPK}^{\rm enr}\)&categoría de estados enriquecidos&
Conserva todas las coordenadas tipadas del núcleo.\\
13&\(\mathcal H\)&relación sobre estados enriquecidos&
Transición Hensel que retiene residuo y cociente.\\
14&\(\mathcal C_{3,1000}\)&estados cruzados
\((N,L,D,K,A,B)\)&
Actualización exacta de cilindros ternario--decimales.\\
15&\(\Sigma_{\rm B}\)&memorias de frontera \(\to\) firma conjunta&
Publica la orientación de bloque conservando su ledger.\\
16&\({\rm EF}\text{--}G_9\)&relación de extensión&
Filtra extensiones compatibles y organiza la genealogía nonádica.\\
17&\(\Gamma_9\)&historias enriquecidas \(\longrightarrow\) historias&
Holonomía de una vuelta de fase con memoria de frontera.\\
18&\(X_\chi\)&\(\varprojlim_m\operatorname{Surv}^{(\chi)}_m\)&
Sección global cuando no vaciedad, compatibilidad y unicidad están
demostradas en toda profundidad.\\
19&\(\mathcal P_{\rm APP}^{(2)}\)&categoría bivariada de caminos&
Transporta simultáneamente la evaluación multiplicativa y la acumulación
aditiva.\\
20&\(F_{\rm obs}\)&\(\mathcal Q_{\rm obs}\to\mathcal Q_{\rm obs}\)&
Cronología observable determinista de dos cursores.\\
21&\(\mathcal A_{42}\)&alfabeto de cuarenta y dos tríadas&
Soporte decimal ampliado de la emisión finita.\\
22&\(\operatorname{Ext}_r\)&
\(\operatorname{Ext}_r\subseteq\mathcal F_q\times\mathcal F_{q'}\)&
Relación tipada de prolongación de fibras.\\
23&\(\mathcal K_{\rm ph}\)&\(\mathcal U_6\times C_9\) sobre sí mismo&
Transferencia de continuaciones compatibles con fase explícita.\\
24&\((P,H,Q)\)&\(\mathcal X_\times\to\mathcal X_\times\)&
Avance, media vuelta y cuarto de giro de semillas.\\
25&\((c_{\rm mode},c_{\rm or})\)&caminos \(\to\mathbb Z^2\)&
Cociclos de cambio de modo y orientación.\\
26&\(E_{20}\)&estados finitos \(\to\) veinte bloques&
Prolongación triádica finita con libro de procedencia.\\
27&\(\mathcal R_{12}\)&historia enriquecida \(\to\) registro de doce
sectores&
Sistema temporal orientado con ruta, firma, acarreo y ledger.\\
28&\(R_{\rm sgn}\)&estado terminal \(\to U_{12}^{\rm sgn}
\subseteq\mathbb Z^{12}\)&
Lector del observable armónico firmado.\\
29&\(\mathcal H_{12}\)&\(\mathbb Z^{12}\to\mathbb Z^{12}\)&
Transformación por tres órbitas de paso tres; inversa
\(\mathcal H_{12}/4\) sobre la imagen pertinente.\\
30&\(\mathcal E_{\rm dod}=(D_3,D_4,Q)\)&
\(\mathbb Z^{12}\to\mathbb Z^{25}\)&
Extractor transversal conjunto de aridades tres y cuatro.\\
31&\(\mathcal C_\alpha\)&bloques y acarreos \(\to\) bloques&
Recurrencia dodecafásica de cierre en base mil.\\
32&\((L_{\rm tick},\chi_{\rm mass})\)&historia enriquecida \(\to\)
carácter temporal&
Interfaz dimensional posterior; no participa en el constructor numérico de
esta obra.\\
33&\(\mathcal W_\pm\)&retícula bidireccional de índice doce&
Lectura en ambos sentidos conservando orientación y procedencia.\\
34&\(\mathcal A_{\rm species}\)&familias de extensión \(\to\) atlas&
Interfaz de realización posterior, no un cuarto primitivo.\\
\bottomrule
\end{longtable}
\endgroup

\section{Tubería causal de actualización}

\begin{definition}[Estado canónico completo]
\label{def:u006f-esquema-estado}
Sea
\[
 \mathcal E:=
 \mathbb Z_{\rm quo}\times\mathbb Z_{\rm car}\times
 \mathsf{Mem}\times\mathsf{Term}\times\mathsf{Cyl}\times
 \mathsf{Fr}\times\mathsf{Pref}\times\Lambda .
\]
El estado enriquecido canónico en el instante \(t\) es
\begin{equation}
 x_t=(p_t,m_t,\sigma_t,b_t,r_t;
 q_t,c_t,\eta_t,s_t,C_t,\partial_t,N_t,\lambda_t)
 \in
 X_9\times\{\Sigma,\Pi\}\times\mathsf{TRIT}\times
 C_{12}\times C_9\times\mathcal E.
 \label{eq:u006f-esquema-estado}
\end{equation}
Las abreviaturas de estado utilizadas en otras unidades son proyecciones de
esta tupla. En particular, no se omiten el cociente \(q_t\), el acarreo
\(c_t\), el terminal \(s_t\), el cilindro \(C_t\), la frontera
\(\partial_t\), el prefijo \(N_t\) ni el libro \(\lambda_t\).
\end{definition}

Sea \(\rho_t\) la flecha observable determinada por \(x_t\). Las leyes de
transporte de U005 y la transición \(\mathcal C_{3,1000}\) determinan un
endomorfismo de memoria. Para
\(z=(q,c,\eta,s,C,\partial,N,\lambda)\in\mathcal E\), se define
\begin{align}
 \mathbf E_t:\mathcal E&\longrightarrow\mathcal E,\notag\\
 \mathbf E_t(z)&=\bigl(
 \mathsf H(\rho_t)q,
 \mathsf C(\rho_t)c+\kappa(\rho_t),
 \mathsf M(\rho_t)\eta,
 \mathsf T(\rho_t)s,
 \mathsf{Cyl}(\rho_t)C,
 \mathsf{Fr}(\rho_t)\partial,
 \mathsf{Pref}(\rho_t)N,
 \Lambda(\rho_t)\lambda
 \bigr).
 \label{eq:u006f-actualizacion-memoria}
\end{align}
Cada componente publicada en \eqref{eq:u006f-actualizacion-memoria} posee el
dominio indicado por su fibra; \(\kappa(\rho_t)\) es la cocadena de acarreo.

Definimos tres endomorfismos del producto canónico. Si
\(\sigma'=M_{\rm ph}(g(t))\) y \(m'=\mu(\sigma',m)\), entonces
\begin{align}
 \mathsf{Sel}_t(p,m,\sigma,b,r;e)
  &:=(p,m',\sigma',b,r;e),\label{eq:u006f-selector-tipado}\\
 \mathsf{Tra}_t(p,m,\sigma,b,r;e)
  &:=(T_{d(t)}p,m,\sigma,b,r;e),\label{eq:u006f-transporte-tipado}\\
 \mathsf{Upd}_t(p,m,\sigma,b,r;e)
  &:=(p,m,\sigma,b',r';\mathbf E_t(e)),
 \label{eq:u006f-memoria-tipada}
\end{align}
donde \((b',r')\) es el destino de \(\rho_t:(b,r)\to(b',r')\). La tubería
causal es la composición de mapas, no la composición de un mapa con un valor
trítico:
\begin{equation}
 \boxed{\mathcal U_t:=
 \mathsf{Upd}_t\circ\mathsf{Tra}_t\circ\mathsf{Sel}_t.}
 \label{eq:u006f-tuberia}
\end{equation}

\begin{theorem}[Cierre de la tubería]
\label{thm:u006f-cierre-tuberia}
La aplicación \(\mathcal U_t\) es un endomorfismo del estado canónico
completo. Los lectores, prolongaciones y proyecciones se aplican después de
\(\mathcal U_t\) sólo cuando su dominio ha sido producido.
\end{theorem}

\begin{proof}
\(\mathsf{Sel}_t\) cambia únicamente modo y TRIT;
\(\mathsf{Tra}_t\) cambia únicamente la posición APP; y
\(\mathsf{Upd}_t\) aplica el producto tipado
\eqref{eq:u006f-actualizacion-memoria} y la flecha de fase. Los tres
codominios son el mismo producto de la
\cref{def:u006f-esquema-estado}; por tanto su composición está definida y
conserva todas las coordenadas del estado.
\end{proof}

\section{Reglas de no identificación}

El inventario impone
\begin{equation}
 \mathcal S_4\neq D_4\neq\operatorname{sd}_4,
 \qquad
 \mathcal R_{12}\neq C_{12}\neq\Theta_{108}\neq\Gamma_{108},
 \qquad
 F_{\rm obs}\neq\mathcal K_{\rm ph},
 \label{eq:u006f-no-identificacion-uno}
\end{equation}
y
\begin{equation}
 U_{12}^{\rm cat}\neq U_{12}^{\rm sgn},
 \qquad
 \mathcal E_{\rm dod}\neq D_{\rm raw}^{90,120}
 \neq\mathcal K_{6\to8}.
 \label{eq:u006f-no-identificacion-dos}
\end{equation}
Son incompatibilidades de tipo, no meras desigualdades de valor.

\begin{criterion}[Auditoría de completitud operatoria]
\label{crit:u006f-completitud}
Una ejecución es operatoriamente completa sólo si:
\begin{enumerate}
 \item declara las entradas de \(\mathfrak O_{\rm TPK}\) empleadas;
 \item registra el dominio y el codominio de cada composición;
 \item conserva el orden causal de \eqref{eq:u006f-tuberia};
 \item distingue dinámica, transferencia, lectura, prolongación y salida;
 \item permite reconstruir desde \(\lambda_t\) cada cambio de bloque y
 cada acarreo;
 \item no promueve una comprobación finita a una sección infinita sin
 probar no vaciedad, compatibilidad y unicidad a toda profundidad.
\end{enumerate}
Cualquier símbolo usado sin tipo o cualquier fusión prohibida por
\eqref{eq:u006f-no-identificacion-uno} y
\eqref{eq:u006f-no-identificacion-dos} constituye un fallo verificable.
\end{criterion}
